\section{System Design and Deployment}
\label{sec:design}

\subsection{The application under test}

We chose a URL-shortener for two reasons. It exercises a realistic mix
of paths --- a write path (POST \texttt{/shorten}) that hashes the
target URL and writes one row into DynamoDB, a read path
(GET \texttt{/r/\{code\}}) that performs a single read-after-write
followed by a 301 redirect, and a CPU-bound path (GET \texttt{/cpu})
that hashes a configurable number of SHA-256 rounds so we can
artificially shift the bottleneck off the network. It is also small
enough that a single-image deployment fits comfortably within the
Learner Lab's 5\,GB ECR quota and short enough to cold-start at all
memory sizes Lambda offers.

The service is implemented in \texttt{FastAPI}~\cite{fastapi} on
Python~3.12. Storage is abstracted behind a small interface so that the
unit tests can use an in-process dictionary and the cloud deployments
use DynamoDB. Eight unit tests, including round-trip POST/GET tests and
a verification of the prefix-stripping middleware that lets us route
\texttt{/ec2/\dots}, \texttt{/fargate/\dots} and \texttt{/lambda/\dots}
to the same application without rewriting URLs at the load balancer,
all pass before deployment.

\subsection{Unified deployment pipeline}

\Cref{fig:arch} shows the topology. A \emph{single} container image
(\texttt{shortlink:v2}) is built once with \texttt{docker buildx} for
\texttt{linux/amd64} and pushed to a shared ECR repository. The same
image is then consumed by:
\begin{itemize}[leftmargin=*,topsep=2pt,itemsep=2pt]
  \item an EC2 Auto Scaling Group of \texttt{t3.small} instances
        running Amazon Linux~2023 with Docker;
  \item an ECS Fargate service running 0.5\,vCPU\,/\,1\,GB tasks;
  \item an AWS Lambda container function with 1\,GB of memory.
\end{itemize}

\begin{figure}[t]
\centering
\begin{tikzpicture}[
  font=\footnotesize,
  node distance=10pt and 22pt,
  box/.style={draw=black!70, thick, rounded corners=2pt, align=center,
              minimum width=70pt, minimum height=28pt, fill=white},
  alb/.style={draw=blue!50!black, very thick, fill=blue!8, rounded corners=2pt,
              minimum width=270pt, minimum height=22pt, align=center, font=\small\bfseries},
  ddb/.style={draw=orange!70!black, very thick, fill=orange!10, rounded corners=2pt,
              minimum width=200pt, minimum height=22pt, align=center, font=\small\bfseries},
  client/.style={draw=black, dashed, rounded corners=2pt, minimum width=80pt,
                 minimum height=22pt, align=center},
  arrow/.style={-{Stealth}, thick},
]
  \node[client] (cl) {client \texttt{t3.small}\\(k6 worker)};
  \node[alb, right=of cl] (alb) {Application Load Balancer};
  \node[box, below=20pt of alb.south west, anchor=north west] (ec2)  {EC2 ASG\\\scriptsize t3.small {\textbullet} x86\_64};
  \node[box, right=of ec2]                                         (fg)   {ECS Fargate\\\scriptsize 0.5\,vCPU {\textbullet} 1\,GB};
  \node[box, right=of fg]                                          (lam)  {Lambda (image)\\\scriptsize 1\,GB {\textbullet} x86\_64};
  \node[ddb, below=20pt of fg] (ddb) {DynamoDB (PAY\_PER\_REQUEST)};

  \draw[arrow] (cl) -- (alb);
  \draw[arrow] (alb.south) ++(-90pt,0) -- node[left,font=\scriptsize]{/ec2/*}     (ec2.north);
  \draw[arrow] (alb.south)              -- node[right,font=\scriptsize]{/fargate/*} (fg.north);
  \draw[arrow] (alb.south) ++( 90pt,0) -- node[right,font=\scriptsize]{/lambda/*} (lam.north);
  \draw[arrow] (ec2.south) -- (ddb.west |- ec2.south) -- (ddb.west);
  \draw[arrow] (fg.south)  -- (ddb.north);
  \draw[arrow] (lam.south) -- (ddb.east |- lam.south) -- (ddb.east);
\end{tikzpicture}
\caption{One container image, three runtimes, one ALB. The same FastAPI
binary serves all three paths; routing is by ALB listener-rule. Storage
is shared so cost differences are confined to the compute layer.}
\label{fig:arch}
\end{figure}

\subsection{From the Learner Lab to a personal AWS account}

The project was originally deployed inside the AWS Academy Learner Lab,
a Vocareum-managed AWS sub-account with two inconvenient properties:
IAM role creation is disabled (our Terraform therefore referenced the
pre-existing \texttt{LabRole} ARN for every service), and public Lambda
Function URLs return \texttt{AccessDeniedException} regardless of
resource-policy configuration. The first pass of measurements was
taken in that environment. Mid-experiment, the lab account was
deactivated by the platform's automated abuse-detection heuristics;
neither Reset nor ticketed intervention by the lab TA could recover it.
All data reported in this paper therefore comes from a fresh personal
AWS account provisioned with the same Terraform, with three
substitutions: (i)~self-managed IAM roles (one EC2 instance profile,
two ECS task roles, one Lambda execution role) replace \texttt{LabRole};
(ii)~Function URLs are dropped --- every paradigm serves through the
shared ALB --- to keep the client-server data path identical across
paradigms and to remove any publicly-reachable invocation vector;
(iii)~Lambda's reserved concurrency is cleared (the new-account
default account-wide limit is 10). A single \texttt{tofu apply}
provisions all 39 resources end-to-end in under seven minutes, and
\texttt{make nuke} returns the account to a clean state including
sweeping stray ENIs or EIPs.
